\documentclass[11pt]{article}

\usepackage[a4paper,margin=2.4cm]{geometry}
\usepackage{amsmath,amssymb,amsthm,mathtools}
\usepackage{array}
\usepackage{booktabs}
\usepackage{longtable}
\usepackage{enumitem}
\usepackage[colorlinks=true,linkcolor=blue,citecolor=blue,urlcolor=blue]{hyperref}
\usepackage{xcolor}
\usepackage[export]{adjustbox}
\usepackage{microtype}
\emergencystretch=2.5em

\title{Typing Rules and Evaluation Rules in Block-MiniJava\\
\large A Formal Reconstruction of the Implemented Static and Dynamic Semantics}
\author{Research note --- Block-MiniJava}
\date{2026-07-20}

% ---------------------------------------------------------------------------
% Inference-rule macro (no external proof-tree package required).
% \irule{RULE-NAME}{premise_1 \quad premise_2 \quad ...}{conclusion}
% ---------------------------------------------------------------------------
\newcommand{\rulename}[1]{\textsc{#1}}
\newcommand{\irule}[3]{%
  \adjustbox{max width=.96\linewidth}{$\displaystyle
    \begin{array}{c}
      #2 \\[2pt] \hline\hline \\[-9pt] #3
    \end{array}
    \;\; \rulename{(#1)}
  $}%
}
\newcommand{\emptyrule}[2]{%
  \adjustbox{max width=.96\linewidth}{$\displaystyle
    \begin{array}{c}
      \hline\hline \\[-9pt] #1
    \end{array}
    \;\; \rulename{(#2)}
  $}%
}

% Judgment shorthands
\newcommand{\Ty}[3]{#1 \vdash #2 : #3}                 % Gamma |- e : tau
\newcommand{\Wf}[2]{#1 \vdash #2 \;\mathsf{ok}}         % Gamma |- s ok
\newcommand{\Sub}[2]{#1 \mathrel{<:} #2}                % tau <: tau'
\newcommand{\Asg}[2]{#1 \mathrel{\rightsquigarrow} #2}  % assignability
\newcommand{\Consist}[2]{#1 \sim #2}                    % gradual consistency
\newcommand{\Step}[1]{\longrightarrow_{\rulename{#1}}}
\newcommand{\mtype}{\mathrm{mtype}}
\newcommand{\fields}{\mathrm{fields}}
\newcommand{\dom}{\mathrm{dom}}
\newcommand{\Hole}{\mathord{?}}

\begin{document}
\maketitle

\begin{abstract}
Block-MiniJava's inspector exposes two kinds of formal artifact at runtime: a
\emph{Typing} tab that renders on-demand natural-deduction derivations for
every method, and three small-step \emph{steppers} (CESK, A~vs~B, Rewrite)
that expose the program's operational semantics one transition at a time.
Both artifacts are produced by the same code that decides whether a program
type-checks and what it evaluates to --- the visualization cannot drift from
the implementation because it \emph{is} the implementation, instrumented to
record its own derivation. This note reconstructs, rule by rule, the static
type system (\S\ref{sec:types}) and the two operational semantics --- an
abstract machine and a substitution-based reduction semantics
(\S\ref{sec:eval}) --- exactly as implemented in
\texttt{src/core/types/typeChecker.ts}, \texttt{src/core/types/classTable.ts},
\texttt{src/core/semantics/minijavaMachine.ts}, and
\texttt{src/core/semantics/minijavaSubstitution.ts}, and traces each rule
family to the programming-language-theory literature it descends from.
\end{abstract}

\tableofcontents

% =============================================================================
\section{Introduction}

Block-MiniJava implements MiniJava (the CSE~401 core, plus a \texttt{String}
extension; see \texttt{MiniJava-Adjusted-BNF.md}) as a block-based language
with:

\begin{itemize}
  \item a three-layer static type system (type grammar, class table with
        nominal subtyping, and a bidirectional checker) that also produces a
        typing \emph{derivation} for the Typing tab, and
  \item two independently implemented, cross-checked operational semantics:
  \begin{enumerate}
    \item \textbf{Vehicle~1} --- a small-step \emph{abstract machine},
          parameterized by a \emph{value model} $\in \{A, B\}$: Model~A is
          faithful MiniJava (a mutable heap, reference semantics, aliasing);
          Model~B is structure-preserving (objects/arrays are inline values,
          field update is functional). Both models share one control-flow
          skeleton, so they can be run in lockstep on the same program (the
          ``A vs B'' tab).
    \item \textbf{Vehicle~2} --- a small-step \emph{substitution} (reduction)
          semantics over the pure fragment of the language: every state is a
          literal block tree, and a step is a literal rewrite of a redex (the
          ``Rewrite'' tab). An \emph{operational correspondence} check
          replays the same program on the Model~B machine and asserts that
          the two traces agree on every \emph{salient} rule.
  \end{enumerate}
\end{itemize}

Every rule below is named exactly as it appears in the source (the
\texttt{rule} field of a \texttt{Derivation}, or the \texttt{lastRule} /
\texttt{rule} field of a machine or substitution step), so this note can be
read side by side with the Typing, CESK, A-vs-B, and Rewrite tabs.

% =============================================================================
\section{The Static Type System}
\label{sec:types}

The type system is implemented in three separable layers, following the
design document's own layering (\texttt{block-minijava-value-and-visualisation-design.md},
\S1):

\begin{enumerate}
  \item \textbf{Type grammar} (\texttt{core/types/ty.ts}) --- the abstract
        syntax of types.
  \item \textbf{Class table} (\texttt{core/types/classTable.ts}) --- name
        resolution, \texttt{extends} chains, nominal subtyping.
  \item \textbf{Checker} (\texttt{core/types/typeChecker.ts}) --- the typing
        and statement well-formedness judgments, which also emit a
        \texttt{Derivation} tree consumed by the Typing tab.
\end{enumerate}

\subsection{Type grammar}

\[
  \tau \;::=\; \mathsf{int} \mid \mathsf{boolean} \mid \mathsf{String}
             \mid \mathsf{int[]} \mid \mathsf{String[]}
             \mid C
             \mid (\tau_1,\ldots,\tau_n)\!\to\!\tau
             \mid \mathsf{Object} \mid \mathsf{Nothing} \mid \Hole
\]

where $C$ ranges over declared class names. $\mathsf{Object}$
(\texttt{Top}) and $\mathsf{Nothing}$ (\texttt{Bottom}) bound the nominal
subtype lattice; the arrow type $(\tau_1,\ldots,\tau_n)\!\to\!\tau$ is
produced only internally, as the result of the $\mtype$ lookup used by
\rulename{T-Invk} (\S\ref{sec:t-invk}) --- there is no arrow-typed surface
syntax. $\Hole$ is a \emph{first-class unfilled-annotation} marker in the
style of Hazel's typed holes (Omar et al.
\cite{omar2017hazelnut,omar2019livelits}): an empty socket, or an expression
whose type could not be determined, types as $\Hole$ and is made consistent
with every other type (\S\ref{sec:assignability}) so that one incomplete
block never cascades into unrelated errors.

\subsection{Contexts, the class table, and subtyping}

A typing context $\Gamma$ is a finite map from identifiers to types. Inside a
method of class $C$, $\Gamma$ is fixed for the whole body:
\[
  \Gamma \;=\; \{\, \mathsf{this} : C \,\} \cup
  \{\, x_i : \tau_i \mid x_i \text{ a parameter} \,\} \cup
  \{\, y_j : \tau_j \mid y_j \text{ a local variable} \,\}
\]
($\mathsf{this}$ is omitted inside \texttt{main}, whose sole binding is
\texttt{args}$:\mathsf{String[]}$). Field lookup is a \emph{separate},
class-indexed environment $\fields(C)$, resolved by walking $C$ and its
ancestors and taking the nearest (most-derived) declaration of a name --- an
identifier that is not in $\Gamma$ but is in $\fields(C)$ is an
\emph{implicit} \texttt{this.}\emph{f} reference. Method lookup is the
analogous $\mtype(m, C)$, giving the $(\tau_1,\ldots,\tau_n)\!\to\!\tau$
signature found by the same nearest-ancestor walk (overriding, not
overloading: MiniJava has one method per name per class). Both
$\fields(\cdot)$ and $\mtype(\cdot,\cdot)$ are named after, and behave
exactly as, the identically named auxiliary functions in Featherweight Java
\cite{igarashi2001fj} (\S\ref{sec:fj-source}).

Subtyping $\Sub{\tau}{\tau'}$ is nominal, seeded from \texttt{extends}:
\[
  \irule{S-Refl}{}{\Sub{\tau}{\tau}}
  \qquad
  \irule{S-Top}{}{\Sub{\tau}{\mathsf{Object}}}
  \qquad
  \irule{S-Bottom}{}{\Sub{\mathsf{Nothing}}{\tau}}
\]
\[
  \irule{S-Extends}{\mathtt{class}\;C\;\mathtt{extends}\;D}{\Sub{C}{D}}
  \qquad
  \irule{S-Trans}{\Sub{C}{D} \quad \Sub{D}{E}}{\Sub{C}{E}}
\]
(the implementation gets \rulename{S-Trans} for free: \texttt{isSubtype}
walks the whole ancestor chain of $C$ looking for $D$, rather than composing
single \texttt{extends} steps.) Two classes with no subtype relationship
except through $\mathsf{Object}$ are otherwise incomparable --- there is no
structural or union/intersection subtyping.

\subsection{Assignability: subtyping plus hole consistency}
\label{sec:assignability}

Every place a value flows into a typed slot (assignment, argument passing,
\texttt{return}) is checked against \emph{assignability}
$\Asg{\tau}{\tau'}$, not raw subtyping:
\[
  \irule{Asg-Sub}{\Sub{\tau}{\tau'}}{\Asg{\tau}{\tau'}}
  \qquad
  \irule{Asg-HoleL}{}{\Asg{\Hole}{\tau'}}
  \qquad
  \irule{Asg-HoleR}{}{\Asg{\tau}{\Hole}}
\]
This is exactly a two-sided \emph{gradual-typing consistency} relation
$\Consist{\tau}{\tau'}$ in the sense of Siek and Taha
\cite{siek2006gradual}, restricted so that the only ``unknown'' type is the
editor's typed hole --- the same design Hazel uses to keep an editable,
possibly-incomplete program meaningfully typed at every keystroke
\cite{omar2017hazelnut}.

\subsection{\texorpdfstring{Expression typing rules ($\mathtt{T}\text{-}\ast$)}{Expression typing rules (T-*)}}

Judgment: $\Ty{\Gamma}{e}{\tau}$ --- ``under $\Gamma$, expression $e$ has
type $\tau$.'' Each rule below is named after the \texttt{rule} string the
checker attaches to the corresponding \texttt{Derivation} node in
\texttt{typeChecker.ts}.

\paragraph{Literals and variables.}
\[
  \emptyrule{\Ty{\Gamma}{n}{\mathsf{int}}}{T-Int}
  \qquad
  \emptyrule{\Ty{\Gamma}{\texttt{"}s\texttt{"}}{\mathsf{String}}}{T-Str}
  \qquad
  \emptyrule{\Ty{\Gamma}{b}{\mathsf{boolean}}}{T-Bool}
  \quad (b \in \{\mathtt{true},\mathtt{false}\})
\]
\[
  \irule{T-Var}{\Gamma(x) = \tau \;\text{ or }\; \fields(C)(x) = \tau}
        {\Ty{\Gamma}{x}{\tau}}
  \qquad
  \irule{T-This}{\text{inside a method of class } C}
        {\Ty{\Gamma}{\mathtt{this}}{C}}
\]
(\rulename{T-Var} is one rule with two side conditions in the code, because
a bare identifier is either a local/parameter lookup in $\Gamma$ or an
implicit field read through $\fields(C)$; the Typing tab's note line reports
which one fired, e.g.\ \texttt{"$\Gamma(v) = \mathsf{int}$"} or
\texttt{"fields(Cell)(f) = int"}. \rulename{T-This} is undefined --- reported
as an error, typed $\Hole$ --- inside \texttt{main}, which has no receiver.)

\paragraph{Object and array construction.}
\[
  \irule{T-New}{C \text{ is declared}, \; C \neq \text{main class}}
        {\Ty{\Gamma}{\mathtt{new}\;C()}{C}}
  \qquad
  \irule{T-NewArr}{\Ty{\Gamma}{e}{\mathsf{int}}}
        {\Ty{\Gamma}{\mathtt{new\ int[}e\mathtt{]}}{\mathsf{int[]}}}
\]

\paragraph{Operators.}
\[
  \irule{T-Arith}{\Ty{\Gamma}{e_1}{\mathsf{int}} \quad \Ty{\Gamma}{e_2}{\mathsf{int}}}
        {\Ty{\Gamma}{e_1 \oplus e_2}{\mathsf{int}}}
  \quad (\oplus \in \{+,-,*,/\})
\]
\[
  \irule{T-Cmp}{\Ty{\Gamma}{e_1}{\mathsf{int}} \quad \Ty{\Gamma}{e_2}{\mathsf{int}}}
        {\Ty{\Gamma}{e_1 \bowtie e_2}{\mathsf{boolean}}}
  \quad (\bowtie \in \{<,\leq,>,\geq\})
\]
\[
  \irule{T-And}{\Ty{\Gamma}{e_1}{\mathsf{boolean}} \quad \Ty{\Gamma}{e_2}{\mathsf{boolean}}}
        {\Ty{\Gamma}{e_1\,\&\&\,e_2}{\mathsf{boolean}}}
  \qquad
  \irule{T-Or}{\Ty{\Gamma}{e_1}{\mathsf{boolean}} \quad \Ty{\Gamma}{e_2}{\mathsf{boolean}}}
        {\Ty{\Gamma}{e_1\,||\,e_2}{\mathsf{boolean}}}
\]
\[
  \irule{T-Not}{\Ty{\Gamma}{e}{\mathsf{boolean}}}
        {\Ty{\Gamma}{!e}{\mathsf{boolean}}}
\]
(Parenthesization, \texttt{mj\_expr\_parens}, is typing-\emph{transparent}:
the checker returns the child's own \texttt{Derivation} unchanged, so no
\rulename{T-Parens} node ever appears in a derivation tree.)

\paragraph{Arrays.}
\[
  \irule{T-Lookup}{\Ty{\Gamma}{e_1}{\mathsf{int[]}} \quad \Ty{\Gamma}{e_2}{\mathsf{int}}}
        {\Ty{\Gamma}{e_1[e_2]}{\mathsf{int}}}
  \qquad
  \irule{T-Length}{\Ty{\Gamma}{e}{\mathsf{int[]}}}
        {\Ty{\Gamma}{e\mathtt{.length}}{\mathsf{int}}}
\]

\paragraph{Strings (extension over the base BNF).}
\[
  \irule{T-CharAt}{\Ty{\Gamma}{e_1}{\mathsf{String}} \quad \Ty{\Gamma}{e_2}{\mathsf{int}}}
        {\Ty{\Gamma}{e_1\mathtt{.charAt(}e_2\mathtt{)}}{\mathsf{String}}}
\]
\[
  \irule{T-Concat}{\Ty{\Gamma}{e_1}{\mathsf{String}} \quad \Ty{\Gamma}{e_2}{\mathsf{String}}}
        {\Ty{\Gamma}{e_1\mathtt{.concat(}e_2\mathtt{)}}{\mathsf{String}}}
  \qquad
  \irule{T-StrLen}{\Ty{\Gamma}{e}{\mathsf{String}}}
        {\Ty{\Gamma}{e\mathtt{.length()}}{\mathsf{int}}}
\]
(\rulename{T-CharAt} returns a \emph{one-character} $\mathsf{String}$, not a
$\mathsf{char}$ --- the language has no character type, matching the design
note's ``strings, no chars'' decision.)

\paragraph{Method invocation.}
\label{sec:t-invk}
\[
  \irule{T-Invk}
    {\Ty{\Gamma}{e_0}{C} \quad
     \mtype(m,C) = (\tau_1,\ldots,\tau_n)\!\to\!\tau \quad
     \Ty{\Gamma}{e_i}{\sigma_i} \quad
     \Asg{\sigma_i}{\tau_i} \;\; (1 \le i \le n)}
    {\Ty{\Gamma}{e_0.m(e_1,\ldots,e_n)}{\tau}}
\]
This is, rule-for-rule, Featherweight Java's \textsc{T-Invk}
\cite[Fig.~1]{igarashi2001fj}, adapted from FJ's purely functional core to
MiniJava's imperative one (arguments are still evaluated for their values,
not for effects that matter to typing). The Typing tab's note line for this
rule literally prints the $\mtype$ judgment, e.g.\
\texttt{"mtype(set, Cell) = (int) -> int"}.

\subsection{\texorpdfstring{Statement well-formedness rules ($\mathtt{WF}\text{-}\ast$)}{Statement well-formedness rules (WF-*)}}

Judgment: $\Wf{\Gamma}{s}$ --- ``under $\Gamma$, statement $s$ is
well-formed.'' Statements have no type of their own (MiniJava's imperative
core is outside Featherweight Java's purely functional scope), so their
judgment concludes $\mathsf{ok}$ rather than a type --- the standard move for
adding an imperative statement layer on top of an expression type system
(cf.\ Pierce \cite[Ch.~13]{pierce2002tapl} on typing references and
commands).

\[
  \irule{WF-Block}{\Wf{\Gamma}{s_1} \;\cdots\; \Wf{\Gamma}{s_k}}
        {\Wf{\Gamma}{\{\,s_1;\ldots;s_k\,\}}}
\]
\[
  \irule{WF-If}
    {\Ty{\Gamma}{e}{\mathsf{boolean}} \quad
     \Wf{\Gamma}{s_1} \;(\text{then}) \quad
     \Wf{\Gamma}{s_2} \;(\text{else})}
    {\Wf{\Gamma}{\mathtt{if}\,(e)\,s_1\,\mathtt{else}\,s_2}}
\]
\[
  \irule{WF-While}
    {\Ty{\Gamma}{e}{\mathsf{boolean}} \quad \Wf{\Gamma}{s}}
    {\Wf{\Gamma}{\mathtt{while}\,(e)\,s}}
\]
\[
  \irule{WF-Print}
    {\Ty{\Gamma}{e}{\mathsf{int}} \;\text{ or }\; \Ty{\Gamma}{e}{\mathsf{String}}}
    {\Wf{\Gamma}{\mathtt{System.out.println(}e\mathtt{)}}}
\]
\[
  \irule{WF-Assign}
    {\Gamma(x) = \tau \;\text{or}\; \fields(C)(x) = \tau \quad
     \Ty{\Gamma}{e}{\sigma} \quad \Asg{\sigma}{\tau}}
    {\Wf{\Gamma}{x = e}}
\]
\[
  \irule{WF-ArrAssign}
    {\Gamma(a) = \mathsf{int[]} \;\text{or}\; \fields(C)(a) = \mathsf{int[]} \quad
     \Ty{\Gamma}{e_1}{\mathsf{int}} \quad \Ty{\Gamma}{e_2}{\mathsf{int}}}
    {\Wf{\Gamma}{a[e_1] = e_2}}
\]

An undeclared variable does not abort the derivation: \rulename{WF-Assign}
and \rulename{T-Var} still produce a node (typed $\Hole$, or with a
\texttt{"}$x \notin \Gamma$\texttt{"} note) so the rest of the program keeps
checking --- the same Hazel-style ``hole, don't cascade'' policy as
\S\ref{sec:assignability}, applied at the statement level.

\subsection{\texorpdfstring{Method well-formedness: $\mathrm{M\text{-}OK}$}{Method well-formedness: M-OK}}

\[
  \irule{M-OK}
    {\Gamma = \{\mathsf{this}\!:\!C, x_1\!:\!\tau_1,\ldots\}
       \quad \Wf{\Gamma}{s_1} \;\cdots\; \Wf{\Gamma}{s_k}
       \quad \Ty{\Gamma}{e_{\mathrm{ret}}}{\sigma}
       \quad \Asg{\sigma}{\tau_{\mathrm{ret}}}}
    {\mathrm{M\text{-}OK} \vdash \mathtt{method}\ C.m\ \mathsf{ok}}
\]
is the derivation \emph{root} the Typing tab renders per method (and once
for \texttt{main}, with no \texttt{return} premise and $\mathsf{this}$
omitted from $\Gamma$). The name and shape mirror Featherweight Java's
method judgment $M~\mathsf{OK}~\mathsf{IN}~C$ \cite[Fig.~1]{igarashi2001fj},
again re-purposed for an imperative statement sequence instead of FJ's
single return expression.

\subsection{Method overriding}

Outside the per-expression/statement rules, the checker enforces one more
well-formedness condition on the class table: an overriding method must have
\emph{identical} parameter types and a \emph{covariant} return type,
\[
  \irule{Override}
    {\Sub{C}{D} \quad \mtype(m,D) = (\tau_1,\ldots,\tau_n)\!\to\!\tau
     \quad \mtype(m,C) = (\tau_1,\ldots,\tau_n)\!\to\!\tau'
     \quad \Sub{\tau'}{\tau}}
    {\text{override of } m \text{ in } C \text{ is well-formed}}
\]
which is FJ's override condition \cite[\S2.3]{igarashi2001fj} generalized
from FJ's invariant-everything rule to Java's actual covariant-return rule
(JLS \S8.4.8.3, since Java~5) --- MiniJava has no overloading, so ``the''
overridden signature is unambiguous.

\subsection{Theoretical sources}
\label{sec:fj-source}

\begin{itemize}
  \item \textbf{Natural-deduction typing judgments, rule-naming convention,
        and the $\mathtt{T}\text{-}\ast$ prefix} follow Pierce's \emph{Types
        and Programming Languages} \cite{pierce2002tapl}, the standard
        textbook presentation of $\Gamma \vdash e : \tau$ typing rules as
        premises-over-a-bar proof trees --- exactly the shape the Typing tab
        renders.
  \item \textbf{The object layer} --- $\mtype$, $\fields$, \rulename{T-Invk},
        \rulename{T-New}, nominal subtyping seeded from \texttt{extends},
        and the $M~\mathsf{OK}$ method judgment --- is Featherweight Java
        \cite{igarashi2001fj}, the standard minimal core calculus for a
        Java-like nominally-typed OO language. FJ's auxiliary functions and
        two of its rule names are used unchanged; MiniJava is close enough
        to FJ's object model (fields, single inheritance, method
        invocation, \texttt{this}, \texttt{new}) that no adaptation was
        needed there.
  \item \textbf{The imperative statement layer} ($\mathtt{WF}\text{-}\ast$)
        has no FJ analogue --- FJ is purely functional, with no assignment,
        loops, or a store. The $\Wf{\Gamma}{s}$ judgment form and its
        composition with the expression layer follow the generic treatment
        of imperative extensions to a typed calculus in Pierce
        \cite[Ch.~13]{pierce2002tapl}.
  \item \textbf{Typed holes and hole consistency} (\S\ref{sec:assignability})
        follow Hazel's live, always-typed editor calculus
        \cite{omar2017hazelnut,omar2019livelits}, itself built on the
        consistency relation from gradual typing \cite{siek2006gradual}.
\end{itemize}

% =============================================================================
\section{Operational Semantics}
\label{sec:eval}

Block-MiniJava implements two independent evaluators over the same block
AST, deliberately keeping them structurally different so that agreement
between them is evidence, not a tautology (design note \S6--\S8;
\texttt{research-notes/abstract-machine-analysis.md}):

\begin{center}
\begin{tabular}{@{}lll@{}}
\toprule
 & \textbf{Vehicle 1 --- abstract machine} & \textbf{Vehicle 2 --- substitution} \\
\midrule
file & \texttt{minijavaMachine.ts} & \texttt{minijavaSubstitution.ts} \\
scope & the full language (Models A and B) & the pure fragment only \\
state & control + frame stack + heap & one block tree \\
step & transition rule on a configuration & literal rewrite of a redex \\
tab & CESK, A vs B & Rewrite \\
\bottomrule
\end{tabular}
\end{center}

\subsection{Vehicle 1: a store-bearing abstract machine, parameterized by value model}

\subsubsection{Configurations}

A machine state is a tuple
\[
  \Sigma \;=\; \langle\, C,\; \varphi\!:\!\Phi,\; H \,\rangle
\]
where $C$ is the \emph{control} (a statement, an expression, a produced
\emph{value}, or \emph{done}); $\varphi\!:\!\Phi$ is a non-empty
\emph{stack of activation frames}, top frame $\varphi$, the rest $\Phi$
(pushed on method call, popped on return); and $H : \mathrm{Loc} \rightharpoonup
\mathrm{HeapObj}$ is the \emph{store}. Each frame is itself a small
CEK-style configuration,
\[
  \varphi \;=\; \langle\, \mathsf{self},\; \rho,\; \kappa \,\rangle,
  \qquad \rho : \mathrm{Var} \rightharpoonup \mathrm{Value}, \quad
  \kappa : \mathrm{Kont}^{*}
\]
with $\rho$ the frame's local environment (parameters and locals, held
\emph{by copy}) and $\kappa$ its pending-work stack, innermost first. This
is a \textbf{CESK machine} (Control, Environment, Store, Kontinuation) in
the sense of Felleisen and Friedman's CEK machine
\cite{felleisen1986control} extended with an explicit, imperative-language
store $S$/$H$ as formalized by Might and Van~Horn's abstracting-abstract-machines
treatment of CESK machines \cite{mightvanhorn2010aam}, further extended with
an explicit \emph{stack of frames} (rather than one shared continuation) to
support method call/return --- the \emph{dump} of Landin's original SECD
machine \cite{landin1964scpp}. The project's own research note classifies
this precisely as ``a CESK-style machine with an essential mutable store,
over an explicit call stack of activation frames''
(\texttt{research-notes/abstract-machine-analysis.md}).

\paragraph{Values diverge by model.}
\[
  v ::= n \mid b \mid \mathrm{s} \mid \mathtt{null}
      \mid \underbrace{\mathrm{Ref}(\ell)}_{\text{Model A only}}
      \mid \underbrace{\mathrm{Obj}(C,\overline{f{=}v}) \mid \mathrm{Arr}(\overline{v})}_{\text{Model B only}}
\]
Under \textbf{Model A}, $\mathrm{Ref}(\ell)$ is the \emph{only} way to reach
an object or array; $H(\ell)$ holds the actual $\mathrm{Obj}/\mathrm{Arr}$,
so aliasing is real and field write is a store mutation. Under
\textbf{Model B}, $H = \varnothing$ always (the ``no-store invariant''):
objects and arrays \emph{are} values, copied by binding, and field/array
update rebuilds a fresh value rather than mutating anything. The two models
share every other part of the machine (control flow, continuations,
frame-stack discipline), so \texttt{step} is one function parameterized by
$\mathrm{model} \in \{A,B\}$, and running both models on the same program
is a genuine lockstep execution, not two separate interpreters that merely
happen to agree.

\subsubsection{Continuations}

\begin{center}
\begin{tabular}{@{}ll@{}}
\toprule
group & continuation frames ($\kappa$) \\
\midrule
statements & $\mathrm{KStmtSeq}, \mathrm{KLoop}, \mathrm{KIf}, \mathrm{KWhile}, \mathrm{KPrint}, \mathrm{KAssign},$ \\
           & $\mathrm{KArrAssignIdx}, \mathrm{KArrAssignVal}$ \\
operators  & $\mathrm{KNot}, \mathrm{KAnd}, \mathrm{KOr}, \mathrm{KBinL}\!\to\!\mathrm{KBinR}$ \\
arrays/strings & $\mathrm{KLookupArr}\!\to\!\mathrm{KLookupIdx}, \mathrm{KLength}, \mathrm{KNewArr},$ \\
                & $\mathrm{KCharAtStr}\!\to\!\mathrm{KCharAtIdx}, \mathrm{KStrLength}$ \\
method call & $\mathrm{KCallRecv}\!\to\!\mathrm{KCallArgs}$ (receiver, then arguments left-to-right, then push a frame) \\
\bottomrule
\end{tabular}
\end{center}

\subsubsection{Transition rules}

Rules are grouped by category; each is named exactly as \texttt{lastRule}
records it. A step $\Sigma \Step{r} \Sigma'$ reads ``$\Sigma$ steps to
$\Sigma'$ by rule $r$.'' Rules marked \emph{salient} correspond 1:1 to a
rewrite of the substitution semantics (\S\ref{sec:correspondence}); rules
marked \emph{administrative} are bookkeeping the substitution semantics
never names (entering a block, dispatching \texttt{if}/\texttt{while},
pushing an operand continuation, \ldots).

\paragraph{Literals and variables (administrative).}
\[
  \irule{lit}{}{\langle \mathrm{Expr}(\ell), \varphi, \Phi, H\rangle \Step{lit} \langle \mathrm{Value}(v(\ell)), \varphi,\Phi,H\rangle}
\]
\[
  \irule{var}{x \in \dom(\rho)}
        {\langle \mathrm{Expr}(x), \langle \mathsf{self},\rho,\kappa\rangle,\Phi,H\rangle \Step{var}
         \langle \mathrm{Value}(\rho(x)), \langle \mathsf{self},\rho,\kappa\rangle,\Phi,H\rangle}
\]
\[
  \irule{field-read}{x \notin \dom(\rho) \quad H(\mathsf{self}).\mathrm{fields}(x) = v \;\text{(A), or}\; \mathsf{self}.\mathrm{fields}(x)=v \;\text{(B)}}
        {\langle \mathrm{Expr}(x), \varphi,\Phi,H\rangle \Step{field\text{-}read} \langle \mathrm{Value}(v),\varphi,\Phi,H\rangle}
\]
\[
  \irule{this}{}
        {\langle \mathrm{Expr}(\mathtt{this}), \langle \mathsf{self},\rho,\kappa\rangle,\Phi,H\rangle \Step{this}
         \langle \mathrm{Value}(\underbrace{\mathrm{Ref}(\mathsf{self})}_{A} \;\text{or}\; \underbrace{\mathsf{self}}_{B}), \varphi,\Phi,H\rangle}
\]

\paragraph{Control flow (administrative except where noted).}
\[
  \irule{seq / seq\text{-}end}{}
    {\langle \mathrm{Done}, \langle\ldots,\mathrm{KStmtSeq}(s_{\mathrm{next}})\!:\!\kappa\rangle\!:\!\Phi,H\rangle
     \Step{seq}
     \begin{cases}
       \langle \mathrm{Stmt}(s_{\mathrm{next}}), \langle\ldots,\mathrm{KStmtSeq}(s_{\mathrm{next}}')\!:\!\kappa\rangle\!:\!\Phi,H\rangle & s_{\mathrm{next}} \neq \bot \\
       \langle \mathrm{Done},\langle\ldots,\kappa\rangle\!:\!\Phi,H\rangle \;(\rulename{seq\text{-}end}) & s_{\mathrm{next}} = \bot
     \end{cases}}
\]
\[
  \irule{if\text{-}true / if\text{-}false}{b \in \{\mathtt{true},\mathtt{false}\}}
    {\langle \mathrm{Value}(b), \langle\ldots,\mathrm{KIf}(s_{\mathtt{then}},s_{\mathtt{else}})\!:\!\kappa\rangle\!:\!\Phi,H\rangle \Step{if\text{-}true/false}
     \langle \mathrm{enter}(b \,?\, s_{\mathtt{then}} : s_{\mathtt{else}}), \ldots\rangle}
\]
\[
  \irule{while\text{-}enter / while\text{-}exit / loop}{}
    {\langle \mathrm{Value}(b),\langle\ldots,\mathrm{KWhile}(e,s)\!:\!\kappa\rangle\!:\!\Phi,H\rangle \Step{while\text{-}enter/exit}
     \begin{cases}
       \mathrm{enter}(s) \text{ with } \mathrm{KLoop}(e,s)\!:\!\kappa & b=\mathtt{true} \\
       \mathrm{Done} \text{ with } \kappa & b=\mathtt{false}
     \end{cases}}
\]
(\rulename{loop} re-focuses the \texttt{while} statement once its body's
$\mathrm{KLoop}$ continuation is reached, restarting the condition.)
\[
  \irule{assign}{}
    {\langle \mathrm{Value}(v), \langle\ldots,\mathrm{KAssign}(x)\!:\!\kappa\rangle\!:\!\Phi,H\rangle \Step{assign}
     \langle \mathrm{Done}, \langle\ldots,\rho[x\!\mapsto\!v],\kappa\rangle\!:\!\Phi,H\rangle}
\]
\[
  \irule{field\text{-}write}{}
    {\langle \mathrm{Value}(v), \langle \mathsf{self}\!=\!\ell,\ldots,\mathrm{KAssign}(f)\!:\!\kappa\rangle\!:\!\Phi,H\rangle \Step{field\text{-}write}
     \langle \mathrm{Done}, \varphi,\Phi,H[\ell.f\!\mapsto\!v]\rangle}
\]
\emph{(Model A: the store mutates.)}
\[
  \irule{field\text{-}update}{}
    {\langle \mathrm{Value}(v), \langle \mathsf{self}\!=\!o,\ldots,\mathrm{KAssign}(f)\!:\!\kappa\rangle\!:\!\Phi,H\rangle \Step{field\text{-}update}
     \langle \mathrm{Done}, \langle o[f\!\mapsto\!v],\ldots\rangle,\Phi,H\rangle}
\]
\emph{(Model B: a fresh object replaces $o$.)}
The pair \rulename{field-write}/\rulename{field-update} is where Models A
and B provably diverge: the built-in \emph{Aliasing Contrast} example
prints \texttt{4141} under A (a shared heap cell observes both writes) and
\texttt{41} under B (each alias holds an independent copy) after the same
51 steps on both machines --- the machine-level demonstration of the
design note's core distinction (\S2--\S4).
\[
  \irule{println}{}
    {\langle \mathrm{Value}(v),\langle\ldots,\mathrm{KPrint}\!:\!\kappa\rangle\!:\!\Phi,H\rangle \Step{println}
     \langle \mathrm{Done},\varphi,\Phi,H\rangle, \; \mathrm{output} \mathrel{+}= \mathrm{show}(v)}
\]
\[
  \irule{array\text{-}write \,/\, array\text{-}update}{}
    {\langle \mathrm{Value}(v), \langle\ldots,\mathrm{KArrAssignVal}(a,i)\!:\!\kappa\rangle\!:\!\Phi,H\rangle \Step{arr\text{-}write/arr\text{-}update}
     \cdots}
\]
(again bifurcating exactly as field write: Model A mutates $H(\ell)$ in
place \rulename{(array-write)}; Model B rebinds $a$ to a new array value
with index $i$ replaced \rulename{(array-update)}.)

\paragraph{Operators (mixed).}
Left-to-right, call-by-value evaluation is enforced by the
$\mathrm{KBinL}\!\to\!\mathrm{KBinR}$ continuation pair (administrative:
\rulename{binop}, \rulename{binop-right}); the salient fold rules are
\[
  \rulename{add}, \rulename{sub}, \rulename{mul}, \rulename{div},
  \rulename{less}, \rulename{leq}, \rulename{gt}, \rulename{geq},
  \rulename{concat}
\]
e.g.
\[
  \irule{add}{}
    {\langle \mathrm{Value}(n_2), \langle\ldots,\mathrm{KBinR}(\mathrm{add},n_1)\!:\!\kappa\rangle\!:\!\Phi,H\rangle
     \Step{add} \langle \mathrm{Value}(n_1{+}n_2),\varphi,\Phi,H\rangle}
\]
\[
  \irule{div}{n_2 \neq 0}
    {\langle \mathrm{Value}(n_2), \langle\ldots,\mathrm{KBinR}(\mathrm{div},n_1)\!:\!\kappa\rangle\!:\!\Phi,H\rangle
     \Step{div} \langle \mathrm{Value}(\mathrm{trunc}(n_1/n_2)),\varphi,\Phi,H\rangle}
\]
(division truncates toward zero and gets \emph{stuck} on $n_2=0$, matching
Java's \texttt{ArithmeticException} rather than silently producing a value).
Boolean connectives short-circuit via dedicated continuations:
\[
  \irule{and\text{-}short\text{-}circuit}{}
    {\langle \mathrm{Value}(\mathtt{false}), \langle\ldots,\mathrm{KAnd}(e_2)\!:\!\kappa\rangle\!:\!\Phi,H\rangle \Step{and\text{-}sc}
     \langle \mathrm{Value}(\mathtt{false}),\varphi,\Phi,H\rangle}
\]
\[
  \irule{or\text{-}short\text{-}circuit}{}
    {\langle \mathrm{Value}(\mathtt{true}), \langle\ldots,\mathrm{KOr}(e_2)\!:\!\kappa\rangle\!:\!\Phi,H\rangle \Step{or\text{-}sc}
     \langle \mathrm{Value}(\mathtt{true}),\varphi,\Phi,H\rangle}
\]
with \rulename{and}/\rulename{or} firing once the (necessarily evaluated)
right operand comes back, and \rulename{not} negating a boolean directly.

\paragraph{Strings.}
\[
  \irule{char\text{-}at}{0 \le i < |s|}
    {\langle \mathrm{Value}(i), \langle\ldots,\mathrm{KCharAtIdx}(s)\!:\!\kappa\rangle\!:\!\Phi,H\rangle
     \Step{char\text{-}at} \langle \mathrm{Value}(s[i]),\varphi,\Phi,H\rangle}
\]
\[
  \irule{str\text{-}length}{}
    {\cdots \Step{str\text{-}length} \langle \mathrm{Value}(|s|),\ldots\rangle}
\]
($i$ out of $[0,|s|)$ is \emph{stuck}: a string index out of bounds, Java's
\texttt{StringIndexOutOfBoundsException}.)

\paragraph{Arrays.}
\[
  \irule{new\text{-}array}{n \ge 0}
    {\langle \mathrm{Value}(n), \langle\ldots,\mathrm{KNewArr}\!:\!\kappa\rangle\!:\!\Phi,H\rangle
     \Step{new\text{-}array}
     \begin{cases}
       \langle \mathrm{Value}(\mathrm{Ref}(\ell)),\varphi,\Phi,H[\ell\!\mapsto\!\overline{0}^n]\rangle & \text{Model A}\\
       \langle \mathrm{Value}(\mathrm{Arr}(\overline{0}^n)),\varphi,\Phi,H\rangle & \text{Model B}
     \end{cases}}
\]
\[
  \irule{array\text{-}read}{0 \le i < |arr|}
    {\cdots \Step{array\text{-}read} \langle \mathrm{Value}(arr[i]),\ldots\rangle}
  \qquad
  \irule{array\text{-}length}{}
    {\cdots \Step{array\text{-}length} \langle \mathrm{Value}(|arr|),\ldots\rangle}
\]

\paragraph{Objects and method call/return.}
\[
  \irule{new}{C \text{ is declared}, \; C \neq \text{main class}}
    {\langle \mathrm{Expr}(\mathtt{new}\;C()), \varphi,\Phi,H\rangle \Step{new}
     \begin{cases}
       \langle \mathrm{Value}(\mathrm{Ref}(\ell)),\varphi,\Phi,H[\ell\!\mapsto\!\mathrm{Obj}(C,\overline{f{=}0})]\rangle & \text{A} \\
       \langle \mathrm{Value}(\mathrm{Obj}(C,\overline{f{=}0})),\varphi,\Phi,H\rangle & \text{B}
     \end{cases}}
\]
\[
  \irule{call}{
    v_0 \text{ resolves to an object of dynamic class } C \quad
    \mtype(m,C) = (\tau_1,\ldots,\tau_n)\!\to\!\tau \quad |args|=n}
    {\langle \mathrm{Value}(v_n), \langle\ldots,\mathrm{KCallArgs}(v_0,\overline{v})\!:\!\kappa\rangle\!:\!\Phi,H\rangle
     \Step{call}
     \langle \mathrm{Stmt}(\mathrm{body}(m,C)), \varphi_{\mathrm{new}}\!:\!\varphi\!:\!\Phi,H\rangle}
\]
where $\varphi_{\mathrm{new}} = \langle \mathsf{self}\!=\!v_0,\;
[x_i\!\mapsto\!v_i]_{i=1}^{n},\; [\,]\rangle$ is a fresh frame --- dynamic
dispatch: the method looked up by $\mtype$ is resolved from $v_0$'s
\emph{dynamic} class, exactly as in FJ/Java, not from the receiver
expression's static type.
\[
  \irule{return\text{-}focus}{\kappa = [\,], \; \text{frame has a return expression not yet focused}}
    {\langle \mathrm{Done}, \varphi,\Phi,H\rangle \Step{return\text{-}focus}
     \langle \mathrm{Expr}(e_{\mathrm{ret}}),\varphi,\Phi,H\rangle}
\]
\[
  \irule{return}{}
    {\langle \mathrm{Value}(v), \varphi\!:\!\Phi,H\rangle \Step{return}
     \langle \mathrm{Value}(v), \Phi, H\rangle}
  \qquad
  \irule{halt}{|\Phi|=0, \; C=\mathrm{Done}}
    {\Sigma \Step{halt} \mathrm{status} := \mathsf{done}}
\]

\subsection{Vehicle 2: substitution-based reduction semantics (pure fragment)}

The Rewrite tab implements a classical \emph{small-step reduction
semantics}: every state is a literal syntax tree, and a step contracts one
\emph{redex} in place, following Plotkin's structural operational semantics
\cite{plotkin1981sos} as refined into the \emph{reduction-semantics-with-evaluation-contexts}
style of Felleisen and Hieb \cite{felleisenhieb1992}, the style taught and
tooled in \emph{Semantics Engineering with PLT Redex}
\cite{felleisenfindlerflatt2009redex} --- the design note's own vocabulary
(``redex,'' ``salient reduction,'' congruence-by-recursive-descent) matches
that tradition directly.

\subsubsection{Scope: the pure fragment}

Vehicle~2 covers only programs whose every method body is a straight-line
sequence of simple assignments followed by a \texttt{return}
(\texttt{analyzeMethod} in \texttt{minijavaSubstitution.ts}); \texttt{if},
\texttt{while}, \texttt{println}, and array writes make a method
\emph{impure}, and the stepper reports \emph{why} rather than silently
changing semantics. \texttt{main} itself is restricted to exactly one
\texttt{System.out.println(\ldots)}. This mirrors the classic move of
carving out a substitution-only calculus (the $\lambda$-calculus core, or
FJ's purely functional method bodies) from a full imperative language, so
that pure substitution --- no environment, no store --- is even meaningful.

\subsubsection{Values and evaluation contexts}

\[
  v \;::=\; n \mid b \mid \mathrm{s} \mid C\{f_1{=}v_1,\ldots,f_k{=}v_k\}
\]
Congruence (``where does the next redex live'') is implemented as
left-to-right, outermost-in recursive descent that stops at the first
non-value child --- an \emph{implicit} evaluation-context grammar
\[
  E \;::=\; [\,] \mid E \oplus e \mid v \oplus E \mid \ldots \mid E.m(\overline{e}) \mid v.m(\overline{v},E,\overline{e})
\]
(one production per operator/constructor, all call-by-value, left
argument first) that \texttt{stepNode}'s recursive structure encodes
directly rather than naming as data --- the standard \emph{reduction-free}
implementation strategy for evaluation contexts \cite{felleisenhieb1992}.

\subsubsection{Reduction rules}

Fold rules (identical arithmetic/relational/string semantics to Vehicle~1,
by construction --- both call the same truncating-division and
1-character-\texttt{charAt} conventions):
\[
  \irule{add/sub/mul/div}{}{n_1 \oplus n_2 \longrightarrow n_1 \oplus_{\mathbb{Z}} n_2}
  \qquad
  \irule{less/leq/gt/geq}{}{n_1 \bowtie n_2 \longrightarrow n_1 \bowtie_{\mathbb{Z}} n_2}
\]
\[
  \irule{and\text{-}short\text{-}circuit}{}{\mathtt{false}\;\&\&\;e_2 \longrightarrow \mathtt{false}}
  \qquad
  \irule{and}{}{\mathtt{true}\;\&\&\;b \longrightarrow b}
\]
\[
  \irule{or\text{-}short\text{-}circuit}{}{\mathtt{true}\;||\;e_2 \longrightarrow \mathtt{true}}
  \qquad
  \irule{or}{}{\mathtt{false}\;||\;b \longrightarrow b}
\]
\[
  \irule{not}{}{!b \longrightarrow \lnot b}
  \qquad
  \irule{parens}{}{(v) \longrightarrow v}
  \qquad
  \irule{concat}{}{\texttt{"}s_1\texttt{"}.\mathtt{concat(}\texttt{"}s_2\texttt{"}\mathtt{)} \longrightarrow \texttt{"}s_1s_2\texttt{"}}
\]
\[
  \irule{char\text{-}at}{0\le i<|s|}{\texttt{"}s\texttt{"}.\mathtt{charAt(}i\mathtt{)} \longrightarrow \texttt{"}s[i]\texttt{"}}
  \qquad
  \irule{str\text{-}length}{}{\texttt{"}s\texttt{"}.\mathtt{length()} \longrightarrow |s|}
\]
\[
  \irule{new}{C \text{ declared}, \; C \neq \text{main class}}
    {\mathtt{new}\;C() \longrightarrow C\{f_1{=}\mathrm{default}(\tau_1),\ldots,f_k{=}\mathrm{default}(\tau_k)\}}
\]

\subsubsection{The $\beta$ rule: method call}

Method call is this semantics' one true $\beta$-rule, and its side
conditions are the most intricate part of the implementation:
\[
  \irule{call}{
    \begin{array}{l}
      v_0 = C\{\overline{f{=}u}\}, \quad \mathrm{mdef}(m,C) = (x_1,\ldots,x_n)\{\,y_1{:=}e_1;\ldots;y_k{:=}e_k;\ \mathtt{return}\ e_{\mathrm{ret}}\,\} \\
      \mathrm{mdef}(m,C) \text{ pure (assignments only)}, \quad n = |v_1,\ldots,v_n| \\
      \rho_0 = [x_i \mapsto v_i]_{i=1}^n, \quad \phi_0 = [\overline{f \mapsto u}] \\
      \forall j\!\in\!1..k:\; w_j = \sigma(e_j; \rho_{j-1}, \phi_{j-1}) \text{ is itself a value}, \quad
        (\rho_j,\phi_j) = \mathrm{update}(\rho_{j-1},\phi_{j-1}, y_j\!\mapsto\!w_j)
    \end{array}
  }
  {v_0.m(v_1,\ldots,v_n) \longrightarrow \sigma(e_{\mathrm{ret}}; \rho_k, \phi_k)}
\]
where $\sigma(e;\rho,\phi)$ substitutes each free identifier $x$ by
$\mathrm{copyFresh}(\rho(x))$ (or $\phi(x)$ if $x$ is a field, i.e.\ an
implicit \texttt{this.}$x$) and each occurrence of \texttt{this} by the
\emph{freshly rebuilt} object $C\{\overline{f \mapsto \phi(f)}\}$ --- so
\texttt{this} inside assignment $j{+}1$ already reflects assignment $j$'s
functional update, matching Model~B's semantics exactly. \textbf{Every
substituted occurrence gets a fresh id}
(\texttt{copyFresh}); this is the \emph{structure-preservation invariant} ---
no two parents ever share a child node, i.e.\ substitution here is the
classical \emph{capture-avoiding, sharing-free} substitution
$e[v/x]$ of the $\lambda$-calculus, not the graph-sharing substitution that
would secretly reintroduce Model~A's aliasing.

The requirement that each $w_j$ \emph{already be a value} after
substitution (checked by \texttt{isValueNode}, else the method is reported
outside the pure fragment) is a deliberate restriction to
\emph{non-duplicating}, call-by-value $\beta$: because $\sigma$ copies
$\rho(x)$ verbatim into every occurrence of $x$, allowing $w_j$ to be an
unevaluated redex would risk re-doing the same computation at every
occurrence --- the classical call-by-name/call-by-value duplication hazard
that $\beta_v$-reduction (call-by-value $\beta$, e.g.\ Plotkin
\cite{plotkin1975callccv}) exists to avoid.

\subsection{Operational correspondence between the two vehicles}
\label{sec:correspondence}

The Rewrite tab's companion check (implemented in \texttt{test/run\_subst.js}
and surfaced live in the UI) replays the same program on the Model~B
machine and asserts that the two traces agree on every \emph{salient} rule
and on the final value, where salient is exactly the set
\begin{center}
\adjustbox{max width=\linewidth}{$\displaystyle
  \{\rulename{add},\rulename{sub},\rulename{mul},\rulename{div},
    \rulename{less},\rulename{leq},\rulename{gt},\rulename{geq},
    \rulename{not},\rulename{and},\rulename{and\text{-}short\text{-}circuit},
    \rulename{or},\rulename{or\text{-}short\text{-}circuit},
    \rulename{new},\rulename{call},\rulename{concat},
    \rulename{char\text{-}at},\rulename{str\text{-}length}\}
$}
\end{center}
Every other machine rule (\texttt{seq}, \texttt{if-*}, \texttt{while-*},
\texttt{var}, \texttt{field-read}, \texttt{push/pop-frame}, the
\texttt{K*} continuation dispatches, \ldots) is \emph{administrative}: it
exists because the machine represents control and environment explicitly,
where the substitution semantics simply rewrites the tree and looks nothing
up. The claim being tested --- that ``environment lookup on the machine is
lazy substitution'' (README, \emph{Semantics \& steppers}) --- is precisely
the informal statement proved formally, for general reduction
semantics/abstract-machine pairs, by Danvy's \emph{syntactic
correspondence} technique (deriving an abstract machine from a reduction
semantics by refocusing, and vice versa)
\cite{danvy2004rational,ageefsen2003refocusing}, and is the operational
half of the progress/preservation methodology that soundness proofs use to
relate a reduction relation to a machine
\cite{wrightfelleisen1994syntactic}. Here the correspondence is checked
\emph{empirically}, by exhaustive comparison of salient-rule traces on a
fixed program suite (\texttt{test/run\_subst.js}, 35 programs at the time of
writing) plus a 400-program differential fuzz suite that additionally
cross-checks both machine models, the substitution stepper, and a legacy
evaluator against each other (\texttt{test/run\_fuzz.js}) --- it is not
machine-checked in a proof assistant, unlike the sibling Live-Types project's
Coq-mechanized metatheory.

\subsection{An auxiliary evaluation rule system: mark-and-sweep garbage collection}

Model~A's heap accumulates unreachable objects as a normal side effect of
`new`/field-write; \texttt{src/core/semantics/gc.ts} implements a classical
two-phase \emph{tracing} collector over that heap:

\begin{enumerate}
  \item \textbf{Mark} --- compute $\mathrm{reachable}(H) \subseteq \dom(H)$
        as the least fixed point of a root set (every $\mathrm{Ref}$
        reachable from a live frame's $\rho$, $\mathsf{self}$, or $\kappa$)
        closed under following object fields and array elements.
  \item \textbf{Sweep} --- restrict $H$ to $\mathrm{reachable}(H)$, dropping
        every other binding.
\end{enumerate}
This is the original mark-and-sweep algorithm as introduced for Lisp by
McCarthy \cite{mccarthy1960recursive} and given its modern, general
treatment (root sets, tracing vs.\ reference counting, generational
variants) in Jones, Hosking, and Moss's \emph{The Garbage Collection
Handbook} \cite{jgh2011gchandbook}. It is deliberately \emph{not} part of
the machine's core transition relation $\Step{}$: GC is triggered
externally (a UI action or an explicit test call to \texttt{runGC}),
observably shrinks $H$, and is proved --- by the GC test suite
(\texttt{test/run\_gc.js}, 14 cases) --- to never change a program's
observable output, i.e.\ it is a \emph{sound} store transformation:
$\forall \Sigma.\ \mathrm{run}(\Sigma) = \mathrm{run}(\mathrm{gc}(\Sigma))$
on output, even though the two runs' heaps differ.

\subsection{Theoretical sources}

\begin{itemize}
  \item \textbf{The CESK abstract machine} follows Felleisen and Friedman's
        CEK machine \cite{felleisen1986control}, extended with an explicit
        store as formalized in Might and Van~Horn's CESK presentation
        \cite{mightvanhorn2010aam}; the explicit call-stack-of-frames
        (rather than one shared continuation) traces back to the
        \emph{dump} component of Landin's SECD machine
        \cite{landin1964scpp}, the historical origin of explicit
        function-call/return machinery in an abstract machine.
  \item \textbf{The substitution/reduction semantics} follows Plotkin's
        structural operational semantics \cite{plotkin1981sos} in the
        evaluation-contexts style of Felleisen and Hieb
        \cite{felleisenhieb1992}, as taught and tooled by \emph{PLT Redex}
        \cite{felleisenfindlerflatt2009redex}; call-by-value $\beta$-reduction
        specifically follows Plotkin's call-by-value $\lambda$-calculus
        \cite{plotkin1975callccv}.
  \item \textbf{The operational-correspondence check} between the two
        vehicles instantiates, empirically rather than by proof, the
        syntactic-correspondence methodology relating reduction semantics
        to abstract machines \cite{danvy2004rational,ageefsen2003refocusing},
        which is also the technique underlying Wright and Felleisen's
        syntactic type-soundness proofs \cite{wrightfelleisen1994syntactic}.
  \item \textbf{Mark-and-sweep garbage collection} follows McCarthy's
        original algorithm \cite{mccarthy1960recursive}, per the modern
        treatment in \cite{jgh2011gchandbook}.
\end{itemize}

% =============================================================================
\section{Summary: rule families and their theoretical lineage}

\begin{longtable}{@{}p{2.6cm}p{5.4cm}p{6.7cm}@{}}
\toprule
\textbf{Rule family} & \textbf{Where} & \textbf{Theoretical source} \\
\midrule
\endhead
$\Gamma \vdash e:\tau$, $\mathtt{T}\text{-}\ast$ &
  \texttt{typeChecker.ts}, expression rules &
  Pierce, \emph{TAPL} \cite{pierce2002tapl} (judgment form and rule-naming
  convention) \\
\rulename{T-Invk}, \rulename{T-New}, $\mtype$, $\fields$, \rulename{M-OK},
  \rulename{Override} &
  \texttt{typeChecker.ts}/\texttt{classTable.ts} &
  Featherweight Java \cite{igarashi2001fj} (names, auxiliary functions, and
  rule shapes reused directly) \\
$\mathtt{WF}\text{-}\ast$, $\Gamma \vdash s\ \mathsf{ok}$ &
  \texttt{typeChecker.ts}, statement rules &
  Imperative-extension treatment of a typed calculus \cite[Ch.~13]{pierce2002tapl};
  no FJ analogue (FJ has no statements) \\
$\Hole$, $\Asg{}{}$ (hole consistency) &
  \texttt{ty.ts}, \texttt{classTable.ts} &
  Hazel's typed holes \cite{omar2017hazelnut,omar2019livelits}, built on
  gradual-typing consistency \cite{siek2006gradual} \\
CESK machine, $\Sigma \longrightarrow \Sigma'$ &
  \texttt{minijavaMachine.ts} &
  Felleisen--Friedman CEK \cite{felleisen1986control}; CESK per
  Might--Van~Horn \cite{mightvanhorn2010aam}; frame stack per Landin's SECD
  dump \cite{landin1964scpp} \\
Substitution rewrite, $e \longrightarrow e'$ &
  \texttt{minijavaSubstitution.ts} &
  Plotkin's SOS \cite{plotkin1981sos}; Felleisen--Hieb evaluation contexts
  \cite{felleisenhieb1992}; PLT Redex style \cite{felleisenfindlerflatt2009redex} \\
\rulename{call} ($\beta$-reduction for method invocation) &
  \texttt{minijavaSubstitution.ts}, \texttt{beta()} &
  Call-by-value $\beta_v$ \cite{plotkin1975callccv} \\
Operational correspondence (salient-rule agreement) &
  \texttt{test/run\_subst.js} &
  Syntactic correspondence between reduction semantics and abstract
  machines \cite{danvy2004rational,ageefsen2003refocusing};
  Wright--Felleisen soundness methodology \cite{wrightfelleisen1994syntactic} \\
Mark-and-sweep GC &
  \texttt{gc.ts} &
  McCarthy \cite{mccarthy1960recursive}; modern treatment in
  \cite{jgh2011gchandbook} \\
\bottomrule
\end{longtable}

% =============================================================================
\section*{Note on scope}
\addcontentsline{toc}{section}{Note on scope}

This note reconstructs the rules \emph{as implemented}; it does not attempt
a soundness proof (progress/preservation) relating \S\ref{sec:types} to
\S\ref{sec:eval}, and the operational-correspondence claim of
\S\ref{sec:correspondence} is currently checked by exhaustive testing, not
mechanized. Both are natural next artifacts if this project follows the
sibling Live-Types work toward a machine-checked metatheory.

% =============================================================================
\begin{thebibliography}{99}

\bibitem{pierce2002tapl}
B.~C. Pierce.
\newblock \emph{Types and Programming Languages}.
\newblock MIT Press, 2002.

\bibitem{igarashi2001fj}
A.~Igarashi, B.~C. Pierce, and P.~Wadler.
\newblock Featherweight Java: a minimal core calculus for Java and GJ.
\newblock \emph{ACM Transactions on Programming Languages and Systems (TOPLAS)}, 23(3):396--450, 2001.
\newblock (Originally OOPSLA 1999.)

\bibitem{omar2017hazelnut}
C.~Omar, I.~Voysey, M.~Hilton, J.~Aldrich, and M.~A. Hammer.
\newblock Hazelnut: a bidirectionally typed structure editor calculus.
\newblock In \emph{POPL}, 2017.

\bibitem{omar2019livelits}
C.~Omar, D.~Moon, A.~Blinn, I.~Voysey, N.~Collins, and R.~Chugh.
\newblock Live functional programming with typed holes.
\newblock In \emph{POPL}, 2019.

\bibitem{siek2006gradual}
J.~G. Siek and W.~Taha.
\newblock Gradual typing for functional languages.
\newblock In \emph{Scheme and Functional Programming Workshop}, 2006.

\bibitem{felleisen1986control}
M.~Felleisen and D.~P. Friedman.
\newblock Control operators, the SECD machine, and the $\lambda$-calculus.
\newblock In \emph{Formal Description of Programming Concepts III}, 1986.

\bibitem{mightvanhorn2010aam}
M.~Might and D.~Van~Horn.
\newblock Abstracting abstract machines.
\newblock In \emph{ICFP}, 2010.

\bibitem{landin1964scpp}
P.~J. Landin.
\newblock The mechanical evaluation of expressions.
\newblock \emph{The Computer Journal}, 6(4):308--320, 1964.

\bibitem{plotkin1981sos}
G.~D. Plotkin.
\newblock A structural approach to operational semantics.
\newblock Technical Report DAIMI FN-19, Aarhus University, 1981.

\bibitem{felleisenhieb1992}
M.~Felleisen and R.~Hieb.
\newblock The revised report on the syntactic theories of sequential control and state.
\newblock \emph{Theoretical Computer Science}, 103(2):235--271, 1992.

\bibitem{felleisenfindlerflatt2009redex}
M.~Felleisen, R.~B. Findler, and M.~Flatt.
\newblock \emph{Semantics Engineering with PLT Redex}.
\newblock MIT Press, 2009.

\bibitem{plotkin1975callccv}
G.~D. Plotkin.
\newblock Call-by-name, call-by-value and the $\lambda$-calculus.
\newblock \emph{Theoretical Computer Science}, 1(2):125--159, 1975.

\bibitem{danvy2004rational}
O.~Danvy.
\newblock A rational deconstruction of Landin's SECD machine.
\newblock In \emph{IFL}, 2004.

\bibitem{ageefsen2003refocusing}
M.~S. Ager, D.~Biernacki, O.~Danvy, and J.~Midtgaard.
\newblock A functional correspondence between evaluators and abstract machines.
\newblock In \emph{PPDP}, 2003.

\bibitem{wrightfelleisen1994syntactic}
A.~K. Wright and M.~Felleisen.
\newblock A syntactic approach to type soundness.
\newblock \emph{Information and Computation}, 115(1):38--94, 1994.

\bibitem{mccarthy1960recursive}
J.~McCarthy.
\newblock Recursive functions of symbolic expressions and their computation by machine, Part I.
\newblock \emph{Communications of the ACM}, 3(4):184--195, 1960.

\bibitem{jgh2011gchandbook}
R.~Jones, A.~Hosking, and E.~Moss.
\newblock \emph{The Garbage Collection Handbook: The Art of Automatic Memory Management}.
\newblock CRC Press, 2011.

\end{thebibliography}

\end{document}
